%=============================================================================
\section{Experiments}
\label{sec:experiments}

We evaluate SART on three synthetic reasoning benchmarks with controlled shortcut injection.

%-----------------------------------------------------------------------------
\subsection{Setup}
\label{sec:exp_setup}

\paragraph{Model and Training.}
We use a GPT-style transformer (2 layers, $d_{\mathrm{model}}=128$, 4 heads, $d_{\mathrm{ff}}=256$, $\approx$277K parameters). All models are trained for 30 epochs with AdamW~\cite{loshchilov2019decoupled}, learning rate $3\times10^{-3}$, weight decay $10^{-5}$, batch size 32, and cosine annealing. For SART: $\alpha=\beta=1.0$, $\tau_A=0.3$, $\tau_R=0.5$, $\lambda=3.0$, $\gamma=1.0$, $\rho=0.5$, selected via grid search over 80 configurations.

\paragraph{Datasets.}
We construct three synthetic datasets with controlled shortcut injection, each with 500 train / 200 val / 300 test samples. Training labels are 70\% shortcut-rule and 30\% true-rule; validation and test always follow the true rule. The test set is split evenly into \emph{clean} and \emph{perturbed} (shortcut contradicts the correct answer).
\begin{itemize}[leftmargin=*, itemsep=1pt]
    \item \textbf{Math-Arithmetic}: Classify $a+b \geq 10$. Shortcut: $a \geq 5$; true rule: evaluate the full sum.
    \item \textbf{Financial-Analysis}: Determine regulatory compliance from four features. Shortcut: revenue $\geq 5$; true rule: margin $\geq 5$ and debt $< 5$.
    \item \textbf{Causal-Reasoning}: Determine causal direction from observed variables and a confounder. Shortcut: correlation $\geq 5$; true rule: $x \geq 5$ and $z < 3$.
\end{itemize}

\paragraph{Baselines.}
We compare SART against eight baselines spanning data-centric, loss-based, distributionally robust, and invariant learning approaches:
(1)~\textbf{Standard Fine-Tuning (SFT)}: vanilla cross-entropy, no shortcut mitigation;
(2)~\textbf{Self-Consistency Decoding}~\cite{wang2023selfconsistency}: majority vote over 5 samples at temperature 0.8;
(3)~\textbf{Data Filtering}: remove samples with answer-token confidence $> 0.90$ after a 3-epoch warmup, then retrain;
(4)~\textbf{JTT (Just Train Twice)}~\cite{liu2021just}: identify error-prone samples in a warmup phase and upweight them during retraining;
(5)~\textbf{Focal Loss}~\cite{lin2017focal}: down-weight well-classified examples via $(1-p_t)^\gamma$ weighting;
(6)~\textbf{Group DRO}~\cite{sagawa2020distributionally}: minimize worst-group loss using shortcut/non-shortcut group annotations;
(7)~\textbf{IRM (Invariant Risk Minimization)}~\cite{arjovsky2019invariant}: learn representations invariant across synthetic environments;
(8)~\textbf{Influence Filtering}~\cite{koh2017understanding}: remove training samples identified as harmful via influence functions after a warmup phase.

\paragraph{Metrics.}
Clean accuracy, perturbed accuracy (robustness), reasoning consistency (fraction of traces matching ground-truth patterns), and shortcut detection F1 (where applicable).

%-----------------------------------------------------------------------------
\subsection{Main Results}
\label{sec:exp_main}

\begin{table}[t]
\centering
\caption{Main results averaged across three synthetic reasoning datasets. \textbf{Bold}: best per column; \underline{underline}: second best. $^\dagger$Requires group annotations. $^\ddagger$Requires multiple environments.}
\label{tab:main_results}
\small
\begin{tabular}{lcccc}
\toprule
\textbf{Method} & \textbf{Accuracy} ($\uparrow$) & \textbf{Robustness} ($\uparrow$) & \textbf{Reasoning} ($\uparrow$) & \textbf{SC Det.\ F1} ($\uparrow$) \\
\midrule
Standard Fine-Tuning        & 68.2 & 24.9  & 53.8 & -- \\
Self-Consistency Decoding   & 69.1 & 24.7  & 53.8 & -- \\
Data Filtering              & 56.4 & 50.4  & 45.1 & 0.60 \\
JTT (Just Train Twice)      & 70.0 & 24.0  & 55.8 & -- \\
Focal Loss                  & 66.7 & 24.4  & 52.0 & -- \\
Group DRO$^\dagger$         & 73.8 & 40.9  & 60.9 & -- \\
IRM$^\ddagger$              & 74.4 & 40.0  & 61.3 & -- \\
V-REx$^\ddagger$            & 74.2 & 44.2  & 61.6 & -- \\
Influence Filtering         & \underline{74.4} & \underline{54.9}  & \underline{63.6} & -- \\
\midrule
\textbf{SART (Ours)}        & \textbf{86.7} & \textbf{68.9}  & \textbf{79.1} & \textbf{0.72} \\
\bottomrule
\end{tabular}
\end{table}

\begin{table}[t]
\centering
\caption{Per-dataset clean and perturbed accuracy (\%). \textbf{Bold}: best; \underline{underline}: second best.}
\label{tab:per_dataset}
\small
\begin{tabular}{llccc}
\toprule
\textbf{Dataset} & \textbf{Method} & \textbf{Acc.\ Clean} ($\uparrow$) & \textbf{Acc.\ Perturbed} ($\uparrow$) & \textbf{Reasoning} ($\uparrow$) \\
\midrule
\multirow{7}{*}{\textsc{Math}}
    & Standard FT          & 76.7 & 13.3 & 76.7 \\
    & Data Filtering       & 44.0 & 30.7 & 44.0 \\
    & Group DRO            & 84.0 & 45.3 & 84.0 \\
    & IRM                  & 84.7 & 44.7 & 84.7 \\
    & Influence Filtering  & \underline{92.7} & \underline{82.7} & \underline{92.7} \\
    & SART (Ours)          & \textbf{96.0} & \textbf{74.0} & \textbf{96.0} \\
\midrule
\multirow{7}{*}{\textsc{Financial}}
    & Standard FT          & 58.7 & 35.3 & 34.7 \\
    & Data Filtering       & 55.3 & \underline{43.3} & 48.7 \\
    & Group DRO            & 65.3 & 36.0 & 43.3 \\
    & IRM                  & 66.0 & 34.0 & 42.7 \\
    & Influence Filtering  & 51.3 & 19.3 & 44.0 \\
    & SART (Ours)          & \textbf{73.3} & \textbf{45.3} & \textbf{68.7} \\
\midrule
\multirow{7}{*}{\textsc{Causal}}
    & Standard FT          & 69.3 & 26.0 & 50.0 \\
    & Data Filtering       & 70.0 & \underline{77.3} & 42.7 \\
    & Group DRO            & 72.0 & 41.3 & 55.3 \\
    & IRM                  & 72.7 & 41.3 & \underline{56.7} \\
    & Influence Filtering  & \underline{79.3} & 62.7 & 54.0 \\
    & SART (Ours)          & \textbf{90.7} & \textbf{87.3} & \textbf{72.7} \\
\bottomrule
\end{tabular}
\end{table}

Table~\ref{tab:main_results} shows that SART substantially outperforms all baselines, achieving 86.7\% clean accuracy (+18.5\,pp over SFT) and 68.9\% robustness (+44.0\,pp), surpassing the strongest baseline (Influence Filtering: 74.4\%/54.9\%) by +12.3\,pp accuracy and +14.0\,pp robustness. Notably, SART achieves these results without requiring group annotations (unlike Group DRO) or multiple training environments (unlike IRM and V-REx), relying solely on gradient-based shortcut detection.

Several observations are noteworthy.
First, Influence Filtering achieves the highest robustness among baselines (54.9\%) but exhibits extreme variance across tasks: it dominates Math-Arithmetic (92.7\%/82.7\%) while collapsing to the worst performance on Financial-Analysis (51.3\%/19.3\%). This brittleness arises because influence-function-based sample removal is effective for single-feature shortcuts but counterproductive when shortcuts involve complex multi-feature interactions.
Second, Data Filtering follows a similar pattern---trading significant accuracy (56.4\%) for robustness (50.4\%), with a catastrophic accuracy drop on Math (44.0\%).
Third, JTT, Focal Loss, and Self-Consistency cluster near the SFT baseline (66--70\% accuracy, 24--25\% robustness), confirming that generic loss modifications and inference-time diversity alone cannot address systematic shortcut injection.

Per-dataset results (Table~\ref{tab:per_dataset}) reveal that \emph{SART achieves the best accuracy, robustness, and reasoning on all three datasets} except Math-Arithmetic robustness (74.0\% vs.\ Influence Filtering's 82.7\%) and Causal-Reasoning robustness where Data Filtering reaches 77.3\% at the cost of severe accuracy degradation. On the most challenging Financial-Analysis task, SART outperforms the next-best method by +7.3\,pp accuracy (vs.\ IRM at 66.0\%) and +2.0\,pp robustness (vs.\ Data Filtering at 43.3\%), while achieving dramatically higher reasoning consistency (68.7\% vs.\ 48.7\%). On Causal-Reasoning, SART achieves 90.7\%/87.3\% accuracy/robustness, far exceeding all baselines.

%-----------------------------------------------------------------------------
\subsection{Ablation Study}
\label{sec:exp_ablation}

\begin{table}[t]
\centering
\caption{Ablation of SART components averaged across three datasets. Grad.\ Align.\ measures cosine similarity between training and validation gradients (higher is better).}
\label{tab:ablation}
\small
\begin{tabular}{lcccc}
\toprule
\textbf{Configuration} & \textbf{Accuracy} ($\uparrow$) & \textbf{Robustness} ($\uparrow$) & \textbf{Grad.\ Align.} ($\uparrow$) & \textbf{SC Det.\ F1} \\
\midrule
SFT (Baseline)              & 68.2 & 24.9  & $-0.18$ & -- \\
\midrule
Reweighting Only            & 76.7 & 51.3  & $-0.11$ & 0.87 \\
Gradient Surgery Only       & 78.2 & 36.0  & $-0.06$ & 0.81 \\
\textbf{SART (Full)}        & \textbf{86.7} & \textbf{68.9}  & $\mathbf{-0.01}$ & 0.72 \\
\bottomrule
\end{tabular}
\end{table}

We isolate each component's contribution (Table~\ref{tab:ablation}). Both components individually provide substantial improvements over SFT: Reweighting provides +8.5\,pp accuracy and +26.4\,pp robustness; Gradient Surgery provides +10.0\,pp accuracy and +11.1\,pp robustness. Critically, the \emph{full combination achieves super-additive gains}: +18.5\,pp accuracy over SFT versus the sum of individual improvements (+18.5\,pp), and +44.0\,pp robustness versus the sum (+37.5\,pp), demonstrating strong synergy between sample-level reweighting and gradient-level correction.

The full method achieves near-zero gradient alignment ($-0.01$), a dramatic improvement over SFT ($-0.18$), indicating nearly complete correction of shortcut-correlated gradient components. Removing Reweighting causes a $-8.4$\,pp accuracy drop and $-32.9$\,pp robustness drop; removing Gradient Surgery causes a $-10.0$\,pp accuracy drop and $-17.6$\,pp robustness drop. Reweighting has the larger impact on robustness, while Gradient Surgery has the larger impact on accuracy. Both components are essential for SART's strong performance.

%-----------------------------------------------------------------------------
\subsection{Hyperparameter Sensitivity}
\label{sec:exp_sensitivity}

\begin{table}[t]
\centering
\caption{Sensitivity of SART to key hyperparameters. Grid search over 80 configurations ($\lambda \times \gamma \times \rho$). Shown: best combined score per parameter value. Combined $= 0.4 \times \text{Acc} + 0.6 \times \text{Rob}$.}
\label{tab:sensitivity}
\small
\begin{tabular}{llccc}
\toprule
\textbf{Parameter} & \textbf{Value} & \textbf{Accuracy} ($\uparrow$) & \textbf{Robustness} ($\uparrow$) & \textbf{Combined} ($\uparrow$) \\
\midrule
\multirow{5}{*}{$\lambda$ (reweighting)}
    & 1.0 & 75.8 & 44.0 & 56.7 \\
    & 1.5 & 81.1 & 52.9 & 64.2 \\
    & 2.0 & 86.9 & 69.6 & 76.5 \\
    & 3.0 & \textbf{87.6} & \textbf{74.2} & \textbf{79.6} \\
    & 5.0 & 86.4 & 73.6 & 78.7 \\
\midrule
\multirow{4}{*}{$\gamma$ (gradient surgery)}
    & 0.3 & 84.0 & 68.4 & 74.7 \\
    & 0.5 & 84.2 & 70.2 & 75.8 \\
    & 0.8 & 82.2 & 56.4 & 66.8 \\
    & 1.0 & \textbf{87.6} & \textbf{74.2} & \textbf{79.6} \\
\midrule
\multirow{4}{*}{$\rho$ (answer suppression)}
    & 0.1 & 86.4 & 73.6 & 78.7 \\
    & 0.3 & 88.0 & 68.0 & 76.0 \\
    & 0.5 & \textbf{87.6} & \textbf{74.2} & \textbf{79.6} \\
    & 0.7 & 86.9 & 69.6 & 76.5 \\
\bottomrule
\end{tabular}
\end{table}

We conduct an exhaustive grid search over 80 configurations evaluating all three datasets per configuration (Table~\ref{tab:sensitivity}). Three key findings emerge:

\textbf{(1) Reweighting strength $\lambda$ is the most influential parameter.} Increasing $\lambda$ from 1.0 to 3.0 yields +11.8\,pp accuracy and +30.2\,pp robustness, with diminishing returns beyond $\lambda \approx 3.0$. This indicates a ``sweet spot'' where shortcut samples are sufficiently downweighted without excessively distorting the training distribution.

\textbf{(2) Strong gradient surgery ($\gamma=1.0$) maximizes performance.} Full-strength gradient projection outperforms moderate values ($\gamma=0.8$: 56.4\% robustness vs.\ $\gamma=1.0$: 74.2\%), suggesting that complete removal of shortcut-aligned gradient components is preferable to partial correction when combined with appropriate reweighting.

\textbf{(3) Answer suppression $\rho$ exhibits a moderate optimum at $\rho=0.5$.} This balances removing shortcut-correlated answer signals while preserving useful gradient information from the answer pathway.

%-----------------------------------------------------------------------------
\subsection{ShortcutScore Validation}
\label{sec:exp_detection}

\begin{figure}[t]
    \centering
    \includegraphics[width=\textwidth]{results/figure3.png}
    \caption{Empirical validation of \textsc{ShortcutScore}. (a) Score vs.\ true shortcut rate (Pearson $r=0.67$). (b) Gradient alignment distribution: shortcut samples (red) skew negative; non-shortcut (green) cluster near zero. (c) Exponential reweighting with $\lambda=3.0$.}
    \label{fig:analysis}
\end{figure}

\textsc{ShortcutScore} achieves average shortcut detection F1 of 0.72 across three datasets, outperforming confidence-based filtering (F1$=$0.60). This confirms that gradient structure captures shortcut reliance that output confidence cannot detect---shortcut samples are correctly labeled and confidently predicted, making them indistinguishable from genuinely easy samples by probability alone. Figure~\ref{fig:analysis}(a) confirms a positive monotonic correlation between $S(s)$ and true shortcut rate ($r=0.67$). The gradient alignment distribution (Figure~\ref{fig:analysis}b) shows clear separation: shortcut samples skew toward negative alignment, while non-shortcut samples cluster near zero.

\subsection{Computational Cost}
\label{sec:exp_efficiency}

The full pipeline completes in under 15 seconds per dataset on Apple M-series hardware, and under 3 minutes on NVIDIA H100 NVL GPUs at server scale. SART incurs approximately $2.5\times$ training overhead versus SFT, attributable to per-sample gradient computation for \textsc{ShortcutScore} (scoring all training samples) and gradient projection during surgery. Validation gradients are recomputed periodically rather than per-step, keeping overhead manageable.
